\chapter{Cline: A Practitioner's Reference}
\label{app:cline}

\begin{remark}[Learning Objectives]
After working through this appendix, the reader should be able to:
\begin{enumerate}
  \item Install Cline as a VS Code extension, configure an LLM provider API key,
        and start a first task session.
  \item Explain Cline's Plan/Act execution model and use Plan Mode to review
        a proposed plan before any files are modified.
  \item Configure a \texttt{.clinerules} file to encode project-specific
        conventions into every Cline session.
  \item Use Cline's checkpoint system to inspect and roll back AI-generated
        changes to any prior state within a session.
  \item Select an appropriate LLM provider and model tier for a given finance
        research task, balancing accuracy, context length, and API cost.
  \item Identify at least three quantitative-finance research tasks that map
        naturally to Cline workflows, and describe how to structure the
        corresponding prompts and permission settings.
\end{enumerate}
\end{remark}

% ---------------------------------------------------------------
\section{What is Cline}
\label{app:cline:overview}

\begin{context}
Cline is an open-source autonomous AI coding agent delivered as a Visual Studio
Code extension.  It was created by developer Saoud Rizwan and first released in
June 2024 under the name \emph{Claude Dev} as part of Anthropic's ``Build with
Claude'' contest.  In October 2024 the project was renamed \emph{Cline} (version
2.0), incorporated as Cline Bot Inc., and opened to the broader AI community.
By mid-2026, the extension had accumulated over 8 million installs across VS Code
and JetBrains, with contributions from more than 250 open-source contributors.
The full source code is available at \url{https://github.com/cline/cline} under
the Apache 2.0 licence.

The defining characteristic of Cline is its \emph{model-agnostic} architecture:
unlike Claude Code, which is tied to Anthropic models, Cline treats the LLM
as a replaceable component and supports more than thirty provider backends,
including Anthropic, OpenAI, Google, AWS Bedrock, Azure OpenAI, DeepSeek, and
local models via Ollama.  The developer supplies their own API key; there is no
bundled model subscription.

For finance researchers, this design has two practical implications.  First, it
makes Cline suitable for environments where model-provider choice is constrained
by institutional compliance policy.  Second, it exposes the per-session token
cost directly: every API call is visible, attributable, and billable, which
encourages token-efficient prompting.
\end{context}

\subsection{Positioning Among AI Coding Assistants}
\label{app:cline:positioning}

Table~\ref{tab:cline-vs-claude-code} compares Cline and Claude Code across
the dimensions most relevant to academic finance research.

\begin{table}[ht]
\centering
\caption{Cline versus Claude Code: key dimensions for finance research use.}
\label{tab:cline-vs-claude-code}
\begin{tabular}{lll}
\toprule
\textbf{Dimension} & \textbf{Cline} & \textbf{Claude Code} \\
\midrule
Interface          & VS Code extension; JetBrains; CLI & Terminal (CLI), IDE panels \\
LLM support        & 30+ providers (model-agnostic) & Anthropic models only \\
Licence            & Apache 2.0 open source & Proprietary (Anthropic) \\
Pricing            & Free extension; BYOK API usage & \$20/mo Max or API usage \\
Browser automation & Yes (Puppeteer/Computer Use) & No \\
MCP support        & Yes (first-class) & Yes \\
Checkpoints        & Shadow Git repo per tool call & Session-scoped snapshots \\
Codebase indexing  & None (dynamic file reads) & Whole-project read on start \\
\texttt{.rules} file & \texttt{.clinerules} per workspace & \texttt{CLAUDE.md} per project \\
\bottomrule
\end{tabular}
\end{table}

Many practitioners use both tools: Cline for IDE-embedded development with
visual diff review and browser-verified UI work; Claude Code for complex
terminal-native tasks that benefit from its whole-project context and richer
skill system.

% ---------------------------------------------------------------
\section{Installation and Setup}
\label{app:cline:install}

\subsection{Prerequisites}
\label{app:cline:prereqs}

Cline requires:
\begin{itemize}
  \item Visual Studio Code 1.84 or later (download from \url{https://code.visualstudio.com}).
  \item An API key from at least one supported LLM provider.
  \item (Optional) Node.js 18+ if you plan to use the Cline CLI separately.
\end{itemize}

\subsection{Installing the Extension}
\label{app:cline:vscode-install}

Install Cline from the VS Code Marketplace:

\begin{enumerate}
  \item Open VS Code and press \texttt{Ctrl+Shift+X} (Windows/Linux) or
        \texttt{Cmd+Shift+X} (macOS) to open the Extensions panel.
  \item Search for \textbf{Cline} (publisher ID: \texttt{saoudrizwan.claude-dev}).
  \item Click \textbf{Install}.
\end{enumerate}

Alternatively, install from the command line:

\begin{minted}{bash}
code --install-extension saoudrizwan.claude-dev
\end{minted}

After installation, the Cline icon appears in the VS Code activity bar.
Click it to open the Cline panel.

\subsection{Configuring a Provider}
\label{app:cline:provider}

On first launch, Cline prompts for an API provider and key.  Click the settings
gear in the Cline panel to open the provider configuration:

\begin{enumerate}
  \item Select a provider from the dropdown (e.g.\ \textbf{Anthropic}).
  \item Enter your API key (e.g.\ \texttt{sk-ant-\ldots}).
  \item Select the model (e.g.\ \texttt{claude-sonnet-4-6}).
  \item Click \textbf{Save}.
\end{enumerate}

API keys are stored in VS Code's secret storage (OS keychain), not in plain
text on disk.  For team environments, use a \texttt{.env} file and configure
Cline to read it from the project root.

Table~\ref{tab:cline-providers} summarises the most commonly used providers
for finance research workloads.

\begin{table}[ht]
\centering
\caption{Selected LLM providers available in Cline, with notes for finance
research use.}
\label{tab:cline-providers}
\begin{tabular}{llp{6cm}}
\toprule
\textbf{Provider} & \textbf{Models} & \textbf{Notes} \\
\midrule
Anthropic   & Claude Sonnet, Opus, Haiku & Strongest reasoning; 200K token context; BYOK \\
OpenAI      & GPT-4o, o1, o3-mini        & Broad library support; function calling \\
Google      & Gemini 1.5 Pro/Flash       & 1M token context; useful for very long filings \\
DeepSeek    & DeepSeek-V3, R1            & Low cost; strong at code generation \\
Ollama      & Any GGUF model             & Fully local; no data leaves the machine \\
AWS Bedrock & Claude, Llama, Titan       & VPC-isolated; suitable for regulated data \\
\bottomrule
\end{tabular}
\end{table}

% ---------------------------------------------------------------
\section{The Plan/Act Execution Model}
\label{app:cline:plan-act}

\begin{context}
Cline's most important safety feature is its \emph{Plan/Act} distinction.
Before making any change to the file system or running any command, Cline
can operate in Plan Mode, where it reads files and reasons about what it
would do---without doing it.  The developer reviews and approves the plan,
then switches to Act Mode to execute.  This separation is especially
valuable for finance research, where a misunderstood task can corrupt a
dataset or delete months of carefully assembled code.
\end{context}

\subsection{Plan Mode}
\label{app:cline:plan-mode}

Enable Plan Mode by toggling the \textbf{Plan} button in the Cline panel
before submitting a task.  In Plan Mode, Cline:

\begin{itemize}
  \item Reads all files it determines are relevant to the task.
  \item Constructs a step-by-step plan describing every file it intends to
        create, modify, or delete and every command it intends to run.
  \item Presents the plan as a structured document in the chat panel.
  \item Takes no action on the file system.
\end{itemize}

Review the plan carefully.  Common issues to check:
\begin{itemize}
  \item Does the plan modify files outside the expected scope?
  \item Does it propose to run commands that could have side effects
        (deleting data, pushing to remote, installing packages globally)?
  \item Are the proposed file paths consistent with the project structure?
\end{itemize}

Once the plan is satisfactory, click \textbf{Approve} (or toggle to Act Mode
and submit).  Cline executes the plan step by step, pausing for approval
at each tool call if the \textbf{Auto-approve} setting is off.

\subsection{Act Mode and Tool Calls}
\label{app:cline:act-mode}

In Act Mode, Cline issues \emph{tool calls} to interact with the environment.
Each tool call is displayed in the Cline panel before execution, giving the
developer an opportunity to approve or reject it.  The built-in tools are:

\begin{description}
  \item[\texttt{read\_file}] Read the content of a file.  Cline uses this
        extensively to understand the codebase before making changes.

  \item[\texttt{write\_file}] Write a new file or overwrite an existing one.
        Cline shows the proposed content in a diff view; the developer approves
        or edits before the file is written.

  \item[\texttt{replace\_in\_file}] Make targeted edits to an existing file
        using \texttt{<\!SEARCH\!>}/\texttt{<\!REPLACE\!>} blocks, analogous
        to Claude Code's Edit tool.

  \item[\texttt{execute\_command}] Run a shell command.  Output is streamed
        in real time.  Cline iterates on failures by reading error output and
        attempting corrections.

  \item[\texttt{browser\_action}] Launch a headless Chromium browser, navigate
        to a URL, click elements, take screenshots, and read the rendered page.
        This enables automated testing of web-based dashboards and data portals.

  \item[\texttt{use\_mcp\_tool}] Call a tool exposed by any connected MCP server,
        extending Cline's reach to external databases, APIs, and custom tools.
\end{description}

\subsection{Auto-Approve Settings}
\label{app:cline:auto-approve}

In the Cline settings panel, individual tool categories can be set to
auto-approve, reducing friction for trusted operations:

\begin{itemize}
  \item \textbf{Read files}: safe to auto-approve for most projects.
  \item \textbf{Write files}: approve manually during early exploration;
        auto-approve after establishing trust in the agent's file-scope
        discipline.
  \item \textbf{Execute commands}: approve manually unless the command set
        is tightly constrained by a \texttt{.clinerules} allowlist.
  \item \textbf{Browser actions}: approve manually; browser automation
        can interact with authenticated sessions and is high-consequence.
\end{itemize}

% ---------------------------------------------------------------
\section{Project Configuration with \texttt{.clinerules}}
\label{app:cline:clinerules}

The \texttt{.clinerules} file, placed at the project root, is injected into
Cline's system prompt at the start of every session.  It is the Cline
equivalent of Claude Code's \texttt{CLAUDE.md}.  Use it to encode:

\begin{itemize}
  \item Project structure and conventions (file naming, directory layout).
  \item Coding standards (PEP~8 for Python, type hints, docstring format).
  \item Constraints on AI behaviour (``never modify files in \texttt{data/raw/}'',
        ``always run \texttt{pytest} before committing'').
  \item Domain context (``this project uses BibTeX for references'',
        ``financial figures are always in USD millions unless noted'').
\end{itemize}

\begin{minted}{markdown}
# Finance Research Lab — Cline Rules

## Project Structure
- `src/` — production Python modules
- `notebooks/` — Jupyter exploratory notebooks (never import from here)
- `data/raw/` — read-only; never modify or delete
- `data/processed/` — output of pipeline scripts
- `tests/` — pytest suite; maintain >80% coverage

## Coding Standards
- Python 3.11+; use type hints on all public functions
- Format with `black`; lint with `ruff`
- Run `pytest` before any commit

## Constraints
- Do not install packages not already in `requirements.txt`
  without explicit approval
- Do not push to remote without explicit instruction
- Do not write API keys or credentials into any file

## Domain Context
- All monetary amounts in USD millions unless stated otherwise
- EDGAR filings are retrieved via the SEC REST API; respect 10 req/s rate limit
- Citation format is BibTeX authoryear (biblatex)
\end{minted}

% ---------------------------------------------------------------
\section{Checkpoints and Recovery}
\label{app:cline:checkpoints}

Every time Cline writes or modifies a file, it records a \emph{checkpoint}
in an internal shadow Git repository maintained inside the VS Code extension's
data directory.  Checkpoints are independent of the project's own Git repository.

To use checkpoints:

\begin{enumerate}
  \item In the Cline chat history, each tool call that modified a file shows
        a \textbf{View diff} link.  Click it to see exactly what the tool call
        changed.
  \item To revert the project to the state before a specific tool call, click
        the \textbf{Restore checkpoint} button next to that tool call.
  \item All file writes made after that tool call are undone; the project
        returns to exactly the checkpoint state.
\end{enumerate}

\begin{remark}[Checkpoints vs.\ project Git history]
\label{rem:cline-checkpoints-vs-git}
Cline's checkpoints are granular (one per tool call) and scoped to the current
VS Code window.  They are a session-level safety net, not a long-term audit
trail.  After a session, commit the desired changes to the project repository
with a descriptive message that attributes the AI contribution.  The project
Git history is the durable record; checkpoints are disposable scaffolding.
\end{remark}

% ---------------------------------------------------------------
\section{Context Management}
\label{app:cline:context}

Cline does not index the codebase upfront or build a vector store.  Instead,
it reads files dynamically as needed---an explicit architectural choice modelled
on how a senior engineer would approach an unfamiliar repository.  When the
context window fills, Cline truncates the oldest messages.  The developer can
also trigger manual context management:

\begin{description}
  \item[\texttt{/newtask}] Summarises the current conversation into a structured
        handoff note (plan, work completed, files changed, next steps) and starts
        a fresh session with a clean context window.  Use this when a long session
        is approaching the context limit or when switching to a new sub-task.

  \item[\texttt{/smol}] Compresses the chat history by summarising older turns
        without starting a new task.
\end{description}

\begin{remark}[Token costs for finance workloads]
\label{rem:cline-costs}
A typical Cline session on a finance data pipeline project---reading several
Python modules, running tests, and iterating on a bug fix---uses between
50{,}000 and 200{,}000 tokens.  At Claude Sonnet rates (approximately \$3/M
input, \$15/M output as of mid-2026), this corresponds to roughly \$0.50--\$2.00
per session.  Heavy users find flat-rate subscriptions more economical; light
academic users benefit from the pay-as-you-go model.
\end{remark}

% ---------------------------------------------------------------
\section{Slash Commands and Context Injection}
\label{app:cline:commands}

\begin{table}[ht]
\centering
\caption{Commonly used Cline slash commands and context operators.}
\label{tab:cline-commands}
\begin{tabular}{ll}
\toprule
\textbf{Command / Operator} & \textbf{Effect} \\
\midrule
\texttt{/newtask}    & Summarise session and start fresh (preserves handoff note) \\
\texttt{/smol}       & Compress context in place \\
\texttt{/newrule}    & Open the \texttt{.clinerules} file for editing \\
\texttt{/reportbug}  & Open a GitHub issue on the Cline repository \\
\texttt{@\textit{file}} & Include a specific file in the next message context \\
\texttt{@\textit{folder}} & Include all files in a folder \\
\texttt{@url}        & Fetch and include the content of a URL \\
\texttt{@problems}   & Include the current VS Code diagnostics (linter errors) \\
\bottomrule
\end{tabular}
\end{table}

The \texttt{@problems} operator is particularly useful in finance data pipelines:
pasting the current set of type errors or linter warnings as context often
allows Cline to fix them in a single tool call without reading the entire file.

% ---------------------------------------------------------------
\section{MCP Integration}
\label{app:cline:mcp}

Cline supports the Model Context Protocol (MCP), a standard interface for
connecting AI agents to external tools, databases, and APIs.  MCP servers
can be installed from the Cline MCP marketplace (accessible from the settings
panel) or configured manually in \texttt{cline\_mcp\_settings.json}.

Finance-relevant MCP servers include:

\begin{description}
  \item[EDGAR MCP] Queries the SEC EDGAR REST API and returns structured filing
        data---accession numbers, document links, XBRL financial statements---
        directly into the Cline context.

  \item[SQL MCP] Connects to a local or remote database (PostgreSQL, DuckDB,
        SQLite) and executes queries on behalf of the AI agent, enabling
        Cline to explore and transform financial datasets without writing
        data to intermediate files.

  \item[Bloomberg Terminal MCP] (enterprise) Routes Bloomberg BQL queries
        through the MCP interface, providing real-time market data to the
        agent within a session.
\end{description}

% ---------------------------------------------------------------
\section{Finance Research Workflows}
\label{app:cline:finance}

\begin{context}
The following examples illustrate how Cline's features combine into coherent
research workflows.  Each workflow begins with a task description submitted
in the Cline chat panel.
\end{context}

\subsection{EDGAR Data Pipeline}
\label{app:cline:finance-edgar}

\textbf{Task prompt:}
\begin{minted}{text}
Using the SEC EDGAR REST API, build a Python function that:
1. Accepts a list of CIK numbers and a date range
2. Retrieves all 10-K filings for each company in that range
3. Extracts the MD&A section (Item 7) from the HTML filing
4. Returns a pandas DataFrame with columns: cik, filing_date, mda_text
Respect the 10 req/s rate limit. Write unit tests.
\end{minted}

\textbf{Expected Cline actions:} (in Plan Mode)
\begin{enumerate}
  \item Read \texttt{src/edgar.py} (if it exists) to understand existing helpers.
  \item Read \texttt{requirements.txt} to check available libraries.
  \item Propose creating \texttt{src/edgar\_mda.py} with the function and
        \texttt{tests/test\_edgar\_mda.py} with unit tests.
  \item No files modified until plan approved.
\end{enumerate}

\subsection{LLM Sentiment Evaluation}
\label{app:cline:finance-sentiment}

\textbf{Task prompt:}
\begin{minted}{text}
Evaluate three sentiment models on the Financial PhraseBank dataset:
FinBERT (local via transformers), GPT-4o-mini (API), and a zero-shot
Claude Haiku prompt. Report accuracy, macro-F1, and cost per 1,000 sentences.
\end{minted}

\textbf{Cline advantage over a static notebook:}  Cline can interleave code
execution with plan refinement.  If the FinBERT model raises an out-of-memory
error on the test machine, Cline reads the error output and proposes batching
or an alternative model---within the same session, without manual intervention.

\subsection{Automated Replication Check}
\label{app:cline:finance-replication}

\textbf{Task prompt:}
\begin{minted}{text}
Replicate Table 3 of Loughran & McDonald (2011) using the EDGAR full-text
search index. Use the LM sentiment dictionary from their website. Compare
your coefficients to the published values and flag any discrepancy > 10%.
\end{minted}

This workflow benefits from Cline's browser automation: it can navigate the
LM dictionary download page, trigger the download, verify the file integrity,
and proceed with the analysis---all in a single session.

% ---------------------------------------------------------------
\section{Security Considerations}
\label{app:cline:security}

\begin{enumerate}
  \item \textbf{Never store API keys in \texttt{.clinerules}.}  That file is
        committed to the repository.  Use environment variables or VS Code's
        secret storage.

  \item \textbf{Review every \texttt{execute\_command} call before approval.}
        Commands can exfiltrate data, install packages, or push code to remotes.
        Enable manual approval for all shell commands in Cline settings.

  \item \textbf{Use Ollama for sensitive data.}  When working with non-public
        client data or proprietary trading signals, configure Cline to use a
        locally hosted model via Ollama.  No tokens leave the machine.

  \item \textbf{Audit the MCP server list.}  Each MCP server has access to
        the tools it exposes.  Only install MCP servers from trusted sources
        and review their source code before connecting them to sessions that
        handle sensitive financial data.
\end{enumerate}

\begin{exercise}[First Cline session]
\label{ex:cline-b}
\textbf{[B]} Install Cline, configure it with an Anthropic or OpenAI API key,
and create a \texttt{.clinerules} file for a hypothetical finance research project.
Submit the following task in Plan Mode: ``Write a Python function that retrieves
the five most recent 10-K filings for a given ticker from EDGAR and returns a
list of (date, accession\_number) tuples.''  Review the plan and describe in
one paragraph what Cline proposed to do and what you would verify before
approving.
\end{exercise}

\begin{exercise}[Checkpoint recovery]
\label{ex:cline-i}
\textbf{[I]} Start a Cline session in Act Mode with auto-approve enabled for
write operations.  Ask Cline to add type hints to all functions in a Python
module of your choice.  After it completes, identify one type annotation that
is incorrect or imprecise (e.g.\ \texttt{Any} where a more specific type
applies), use the checkpoint viewer to locate the specific tool call that
introduced it, and restore the checkpoint to the state before that tool call.
Then re-submit the task with an explicit instruction to avoid the
imprecision you identified.
\end{exercise}

\begin{exercise}[Multi-provider cost comparison]
\label{ex:cline-a}
\textbf{[A]} Configure Cline with two providers: Claude Sonnet and DeepSeek-V3.
Run identical sessions---same task, same project, same \texttt{.clinerules}---with
each provider.  Record the token counts (input and output) reported by Cline's
cost tracker for each session.  Evaluate the quality of the outputs on three
criteria: correctness of generated code (does it run and pass tests?), adherence
to \texttt{.clinerules} conventions, and appropriateness of the generated
commit message.  Present a cost-quality trade-off analysis and recommend a
provider configuration for a research group with a \$50/month AI budget.
\end{exercise}
